\chapter{Estado del Arte}
\label{cap:estado}

\section{Metodología de la revisión}

\subsection{Objetivos y preguntas de revisión}

La revisión busca identificar y evaluar el estado del arte en técnicas de
aprendizaje profundo aplicables a una solución automatizada que produzca
anotaciones detalladas, que incluyen localización tiempo--frecuencia,
identificación de especie y clasificación del tipo de vocalización, sobre grabaciones de
monitoreo acústico pasivo de primates amazónicos. Se organiza en torno a cinco
preguntas:

\begin{description}
    \item[P1.] ¿Qué estrategias de aprendizaje profundo han obtenido mejores
          resultados en la detección tiempo--frecuencia de vocalizaciones en
          grabaciones de monitoreo acústico pasivo capturadas en entornos ruidosos?

    \item[P2.] ¿Qué estrategias han obtenido mejores resultados en la
          clasificación jerárquica fina de los eventos detectados, identificando de
          forma simultánea la especie y el tipo específico de vocalización?

    \item[P3.] ¿Qué técnicas de procesamiento resultan más efectivas para
          mitigar el ruido ambiental y la variabilidad propia de sensores de bajo
          costo como AudioMoth?

    \item[P4.] ¿Qué conjuntos de datos públicos existen con vocalizaciones
          etiquetadas, qué formatos de anotación emplean, y qué estrategias de aumento
          de datos compensan mejor la escasez y el desbalance de clases?

    \item[P5.] ¿Cómo se adaptan mediante transferencia de aprendizaje los
          modelos preentrenados en audio general o en bioacústica de aves a las tareas
          de detección (P1) y clasificación jerárquica (P2)?
\end{description}

\subsection{Estrategia de búsqueda}

La búsqueda se ejecutó sobre dos bases indexadas. Scopus cubre publicaciones
seriadas de más de cinco mil editores en 140 países e incorpora una herramienta
de búsqueda asistida que agrupa los artículos relacionados con una consulta. Web
of Science aporta herramientas de visualización sobre la relevancia y las
tendencias de las revistas dentro de sus categorías temáticas. La
Tabla~\ref{tab:cadenas} recoge las cadenas empleadas y el número de documentos
recuperado en cada motor.

\begin{table}[htbp]
    \centering
    \small
    \begin{threeparttable}
        \caption[Cadenas de búsqueda]{Cadenas de búsqueda y documentos recuperados por motor}
        \label{tab:cadenas}
        \begin{tabular}{@{}L{2.4cm}L{9.8cm}r@{}}
            \toprule
            \textbf{Motor}                  & \textbf{Cadena de búsqueda} & \thead{Docs.} \\
            \midrule
            Scopus                          &
            \texttt{TITLE-ABS-KEY} (\textit{deep learning} OR \textit{neural
                network*} OR \textit{artificial intelligence} OR
            \textit{transformer*} OR \textit{machine learning}) AND
            \texttt{TITLE-ABS-KEY} (\textit{sound event detection} OR
            \textit{acoustic event detection} OR \textit{audio event
                localization} OR \textit{bioacoustic detection}) AND
            \texttt{TITLE-ABS-KEY} (\textit{biodiversity} OR \textit{wildlife} OR
            \textit{animal vocalizations} OR \textit{primates} OR
            \textit{anthropogenic noise} OR \textit{forest monitoring} OR
            \textit{ecological monitoring}) & 33                                          \\
            \addlinespace
            Web of Science                  &
            \texttt{TS} = (mismos tres bloques de términos, unidos por
            \texttt{AND})                   & 6                                           \\
            \midrule
            \textbf{Total}                  &                             & \textbf{39}   \\
            \bottomrule
        \end{tabular}
        \begin{tablenotes}[flushleft]
            \footnotesize
            \item Los tres bloques combinan técnica, tarea y dominio de
            aplicación. El asterisco es el comodín de truncamiento de ambos
            motores.
        \end{tablenotes}
    \end{threeparttable}
\end{table}

\subsection{Criterios de inclusión y exclusión}

Los resultados se cribaron con los criterios de la Tabla~\ref{tab:criterios}.
Dos de ellos hacen el trabajo principal de filtrado y merecen explicitarse: el
criterio de detección de eventos descarta los trabajos que solo clasifican la
grabación completa sin localizar nada dentro de ella, y el criterio de
clasificación jerárquica descarta los clasificadores binarios o de un solo nivel
taxonómico. Son precisamente los dos ejes en los que se sitúa el vacío descrito
en la sección~\ref{sec:vacio}.

\begin{table}[htbp]
    \centering
    \caption[Criterios de inclusión y exclusión]{Criterios de inclusión y exclusión aplicados a los resultados}
    \label{tab:criterios}
    \small
    \begin{tabular}{@{}L{2.8cm}L{5.6cm}L{5.6cm}@{}}
        \toprule
        \textbf{Criterio} & \textbf{Inclusión}                                         & \textbf{Exclusión}                                        \\
        \midrule
        Temporalidad      & Publicaciones de 2020 en adelante                          & Anteriores a 2020                                         \\
        \addlinespace
        Calidad           & Revistas Q1 o conferencias de alto impacto                 & Revistas no indexadas o de cuartiles inferiores           \\
        \addlinespace
        Enfoque           & Técnicas de aprendizaje profundo                           & Aprendizaje automático clásico sin componente profundo    \\
        \addlinespace
        Análisis de señal & Análisis tiempo--frecuencia                                & Análisis exclusivamente temporal o frecuencial            \\
        \addlinespace
        Dominio           & Bioacústica, con preferencia por mamíferos y primates      & Reconocimiento de voz humana sin transferencia al dominio \\
        \addlinespace
        Entorno           & Entornos naturales ruidosos, preferentemente tropicales    & Entornos exclusivamente de laboratorio                    \\
        \addlinespace
        Detección         & Métodos que localizan el evento en tiempo y frecuencia     & Métodos que clasifican la grabación completa              \\
        \addlinespace
        Clasificación     & Clasificación multinivel de especie y tipo de vocalización & Clasificadores binarios o de un solo nivel                \\
        \addlinespace
        Desbalance        & Técnicas para desbalance de clases y escasez de datos      & Métodos que exigen grandes volúmenes etiquetados          \\
        \addlinespace
        Evaluación        & Métricas completas de precisión, \textit{recall} y F1      & Trabajos sin evaluación cuantitativa                      \\
        \addlinespace
        Idioma            & Inglés o español                                           & Otros idiomas                                             \\
        \bottomrule
    \end{tabular}
\end{table}

\section{Resultados de la revisión}

Ninguno de los trabajos recuperados evalúa arquitecturas de aprendizaje profundo
para la detección y clasificación de vocalizaciones de primates en el contexto
amazónico. Los resultados que siguen se apoyan, por tanto, en la extrapolación
desde dominios bioacústicos adyacentes pero más desarrollados, sobre todo la
detección de aves y la localización de eventos sonoros, que enfrentan
desafíos análogos de ruido ambiental, variabilidad entre sensores y escasez de
datos etiquetados \citep{lostanlen_robust_2019,meedeniya_survey_2023}.

\subsection{P1: Estrategias de detección tiempo--frecuencia}

Las redes convolucionales son la base de los sistemas que operan sobre
espectrogramas. Un modelo convolucional de arquitectura relativamente simple ya
supera de forma clara a métodos tradicionales como el flujo espectral o los
detectores basados en energía, aunque adolece de falta de robustez y de
generalización ante la variabilidad del ruido \citep{lostanlen_robust_2019}. Las
redes convolucionales recurrentes añaden capas recurrentes para modelar
dependencias temporales más largas y fueron la arquitectura predominante en las
ediciones iniciales de los desafíos de localización y detección de eventos
sonoros \citep{politis_dataset_2020}. Las variantes residuales facilitan el
entrenamiento de redes más profundas sin degradación \citep{he_deep_2016}, y son
la base de BirdNET \citep{kahl_birdnet_2021}.

Los transformadores aparecen como innovación arquitectónica en detección
bioacústica de pocos ejemplos, si bien la mayoría de sistemas exitosos de ese
desafío siguieron usando redes convolucionales \citep{nolasco_learning_2023}. En
detección con salida bidimensional, el trabajo relevante sustituye los
\textit{backbones} preentrenados en imágenes por representaciones nativas de
audio y añade una pirámide multiescala, precisamente porque las arquitecturas de
tipo transformador de espectrograma producen representaciones de una sola escala
\citep{zhu_sound_2025}.

Del lado de la representación de entrada, la normalización de energía por canal
fue el contribuyente principal a la mejora de precisión sobre la línea base de
espectrograma log-mel \citep{lostanlen_robust_2019}, y también la representación
mayoritaria entre los sistemas exitosos del desafío de detección con pocos
ejemplos \citep{nolasco_learning_2023}. La elección de la representación
resultó, en esos trabajos, tan determinante como la elección de la arquitectura
\citep{bonet-sola_comparative_2021}.

\subsection{P2: Clasificación jerárquica fina}

La revisión no identificó estrategias diseñadas y evaluadas explícitamente para
la clasificación jerárquica fina que asigna de forma simultánea la especie y el
tipo de vocalización. La literatura recuperada se concentra en tres tareas
vecinas pero distintas: la detección de eventos, que localiza sin etiquetar en
detalle \citep{lostanlen_robust_2019,politis_dataset_2020}; la clasificación de
especie, multiclase o multietiqueta, que identifica qué especies están presentes
BirdNET reconoce 984 especies de aves \citep{kahl_birdnet_2021}, y
\citet{kath_leveraging_2024} entrenan clasificadores lineales sobre
\textit{embeddings} para anuros; y la clasificación de eventos ya detectados
en un conjunto plano de clases predefinidas \citep{politis_dataset_2020}.

La necesidad de distinguir tipos de llamada o sílabas se reconoce como un
desafío de clasificación de grano fino en bioacústica, pero se enuncia como
problema abierto y no como capacidad demostrada \citep{nolasco_learning_2023}.
Este es, en la revisión, el hueco más nítido y el que sostiene el eje de tarea
de la sección~\ref{sec:vacio}.

\subsection{P3: Mitigación de ruido y variabilidad entre sensores}

Ningún trabajo recuperado prueba técnicas específicamente sobre ruido amazónico,
pero varios abordan el ruido ambiental heterogéneo y la variabilidad entre
sensores en despliegues de monitoreo acústico pasivo. Dos resultados son
directamente aplicables.

El primero es la normalización de energía por canal como preprocesamiento del
espectrograma: al combinar compresión de rango dinámico y control adaptativo de
ganancia, suprime el ruido estacionario y realza los eventos transitorios, y con
sus parámetros ajustados a entornos exteriores ruidosos produce distribuciones
de magnitud consistentes y cuasi-gaussianas entre sensores distintos
\citep{lostanlen_robust_2019}. El segundo son las redes adaptativas al contexto,
que procesan en una rama auxiliar estadísticas espectrales resumidas a largo
plazo, mediante percentiles del espectrograma, y adaptan el modelo a las condiciones
de cada punto de monitoreo, mejorando la robustez frente a la variabilidad
espacial del ruido \citep{lostanlen_robust_2019}. Ambas técnicas resultaron
complementarias: su combinación fue lo que produjo detección robusta a escala de
red de sensores.

\subsection{P4: Conjuntos de datos y aumento de datos}

Los conjuntos públicos disponibles se concentran de forma abrumadora en aves:
BirdVox, Xeno-canto y la Macaulay Library. Fuera de ese taxón, la revisión
identificó AnuraSet para anuros, la base de Watkins para mamíferos marinos y los
conjuntos de DCASE para sonidos generales. No se encontró ningún conjunto
público a gran escala etiquetado específicamente para vocalizaciones de
primates, lo que confirma por una segunda vía el eje de dominio de la
sección~\ref{sec:vacio}.

En cuanto a formatos, los eventos localizados en tiempo se distribuyen
típicamente como archivos delimitados con tiempo de inicio, tiempo de fin y
etiqueta de clase \citep{nolasco_learning_2023}; en tareas de localización se
añade la dirección de llegada \citep{politis_dataset_2020}; y en repositorios
grandes como Xeno-canto lo habitual son etiquetas débiles de presencia o
ausencia a nivel de grabación completa \citep{kahl_birdnet_2021}. Ninguno de
esos formatos registra la banda de frecuencia del evento.

Frente a la escasez y al desbalance, las técnicas reportadas son la adición de
ruido de fondo realista, el desplazamiento de tono, el estiramiento temporal, el
enmascaramiento de bloques de tiempo o frecuencia sobre el espectrograma y el
entrenamiento con mezcla de ejemplos
\citep{kahl_birdnet_2021,lostanlen_robust_2019,politis_dataset_2020}. El
sobremuestreo de clases minoritarias es la respuesta habitual al desbalance,
mientras que el aprendizaje sensible al costo fue probado sin mejora apreciable
\citep{kahl_birdnet_2021}. Cuando los ejemplos etiquetados son muy escasos, las
redes prototípicas permiten adaptarse a clases nuevas con apenas cinco ejemplos
\citep{nolasco_learning_2023}.

\subsection{P5: Transferencia de aprendizaje}

La transferencia de aprendizaje es la estrategia central cuando los datos
etiquetados son limitados \citep{kath_leveraging_2024}. El hallazgo más útil de
la revisión es que la proximidad del dominio de origen importa más que el tamaño
del modelo: BirdNET, entrenado en aves, transfiere mejor a la clasificación de
anuros que modelos entrenados en audio general o en imágenes
\citep{kath_leveraging_2024}. Eso lo posiciona como candidato razonable para
transferir a primates.

Sobre el método de adaptación, la revisión distingue dos rutas. La extracción de
características usa el modelo preentrenado como extractor fijo y entrena un
clasificador simple sobre sus \textit{embeddings}; la penúltima capa de BirdNET
resultó particularmente efectiva en esa función \citep{kath_leveraging_2024}. El
ajuste fino reentrena parcial o totalmente los pesos con los pocos datos de la
tarea objetivo, y exige cuidado para no sobreajustar \citep{nolasco_learning_2023}.
Ninguno de los trabajos documenta la adaptación de estos modelos a la detección
con salida bidimensional ni a la clasificación jerárquica de primates.

\rotulo{Detección y robustez (P1, P3).} No se identificó ningún estudio que
aplique y evalúe modelos de aprendizaje profundo a la detección
tiempo--frecuencia de primates amazónicos. La literatura de dominios vecinos sí
ofrece estrategias trasladables: las arquitecturas convolucionales y
convolucionales recurrentes como base
\citep{lostanlen_robust_2019,politis_dataset_2020}, las variantes residuales y
los transformadores para secuencias largas
\citep{he_deep_2016,nolasco_learning_2023}, y, frente a la variabilidad del
ruido, la normalización de energía por canal junto con las redes adaptativas
al contexto, cuya combinación resultó la más efectiva a escala de red de
sensores \citep{lostanlen_robust_2019}.

\rotulo{Clasificación jerárquica (P2).} No se hallaron estrategias explícitamente
diseñadas y evaluadas para asignar de forma conjunta especie y tipo de
vocalización. Existen enfoques de clasificación multietiqueta sobre
\textit{embeddings} preentrenados \citep{kath_leveraging_2024}, pero la tarea
jerárquica específica sigue requiriendo desarrollo propio.

\rotulo{Datos y anotaciones (P4).} Se confirma la ausencia de conjuntos públicos
a gran escala etiquetados para primates y adecuados para aprendizaje profundo:
el esfuerzo de la comunidad se concentra en aves. Los formatos de anotación de
uso común registran tiempo de inicio, tiempo de fin y clase
\citep{nolasco_learning_2023}, y descartan la banda de frecuencia, lo que
significa que incluso los conjuntos disponibles no servirían como referencia
para la tarea de esta tesis sin reanotación.

\rotulo{Transferencia (P5).} La transferencia desde dominios cercanos es la
respuesta más viable a la escasez de datos etiquetados, con evidencia de
transferencia exitosa de aves a otros taxones
\citep{kath_leveraging_2024,nolasco_learning_2023}. Queda abierto cómo aplicarla
cuando la salida deseada no es una etiqueta por segmento sino una caja
tiempo--frecuencia, que es exactamente el problema que aborda el
Capítulo~\ref{cap:deteccion}.

Los dos ejes del vacío descrito en la sección~\ref{sec:vacio} quedan
confirmados de forma independiente por esta
revisión: no hay trabajo publicado sobre primates neotropicales en este
contexto (eje de dominio) y no hay estrategias establecidas para la salida
bidimensional con etiquetado jerárquico (eje de tarea). El cruce de ambos ejes
es la posición que ocupa esta tesis.


% ===================================================================
% CAPÍTULO 4
% ===================================================================
